
% ---------------------------------------------------------------------------
% Case 4: Obermeyer
% ---------------------------------------------------------------------------
\begin{frame}[allowframebreaks]{Case 4: A Health Risk Algorithm and Its Proxy Label (2019)}

    \textbf{The system.} A widely deployed algorithm identified patients who should be enrolled in
    high-intensity care-management programmes. It predicted \textbf{future health-care costs} as a
    stand-in for \textbf{future health needs}.

    \vspace{0.8em}

    \begin{block}{The defect}
        Less money is spent on Black patients at the same level of illness. Predicting cost
        therefore predicts \emph{spending}, not \emph{sickness} --- and the algorithm concluded that
        equally sick Black patients were healthier.
    \end{block}

    \vspace{0.8em}

    \begin{itemize}
        \setlength\itemsep{0.4em}
        \item At a given risk score, Black patients were measurably sicker (more uncontrolled
              chronic conditions).
        \item Reformulating the label removed most of the disparity: the share of Black patients
              flagged for extra help rose from \textbf{17.7\% to 46.5\%}.
        \item The sensitive attribute was \textbf{not} an input. The \textbf{label} carried the bias.
    \end{itemize}

    \vspace{0.4em}
    {\scriptsize Z.~Obermeyer, B.~Powers, C.~Vogeli and S.~Mullainathan, ``Dissecting racial bias in
    an algorithm used to manage the health of populations,'' \emph{Science} 366(6464):447--453, 2019.}
\end{frame}
